\subsubsection{有界算子, 延拓}\label{subsec:bounded-operators-extensions}

现在修改有界算子及其范数的定义, 以适应这里更一般的算子定义.

\begin{Definition}
\label{def:bounded-operator-specified-domain}
若存在常数 $C$, 使
\begin{equation}\label{eq:bounded-domain-operator-norm-inequality}
\|A(x)\|_Y\leq C\|x\|_X\quad\text{for every }x\in\mathcal{D}(A).
\end{equation}
则称从 $X$ 到 $Y$ 的线性算子 $(\mathcal{D}(A),A)$ \textbf{有界}. 若不存在这样的 $C$, 则称算子\textbf{无界}.

有界算子的范数是使 \eqref{eq:bounded-domain-operator-norm-inequality} 成立的最小 $C$, 即非负数
\begin{equation}\label{eq:bounded-domain-operator-norm-definition}
\|A\|:=\sup_{\substack{x\in\mathcal{D}(A)\\\|x\|_X\ne0}}\dfrac{\|A(x)\|_Y}{\|x\|_X}=\sup_{\substack{x\in\mathcal{D}(A)\\\|x\|=1}}\|A(x)\|_Y.
\end{equation}

\end{Definition}

注意, 若 $\mathcal{D}(A)=X$, 则这一定义与 Definition~\ref{def:bounded-linear-mapping} 一致. 连续性等性质的定义容易推广到这个新版本的定义.

以下引理说明了为什么只有在研究无穷维空间时才需要考察无界算子.

\begin{Lemma}
\label{lem:finite-dimensional-domain-operators-bounded}
所有定义域 $\mathcal{D}(A)$ 为有限维的算子都有界.
\end{Lemma}

\begingroup
\postdisplaypenalty=10000
\renewcommand{\endProof}{\par\nobreak\hfill$\blacksquare$\par\vspace{0.1cm}}
\begin{Proof}
设 $\{e_i\}_{i=1}^n$ 为 $\mathcal{D}(A)$ 的一组基. 由 Lemma~\ref{lem:linear-mapping-bounded-continuity-equivalence}, 只需证明每当 $x_j\to0$ 时都有 $A(x_j)\to0$. 为说明这一点, 注意可写成 $x_j=\sum_{i=1}^n\alpha_{i,j}e_i$. 此外, 每当 $\lim_{j\to\infty}x_j=0$ 时, 有
\begin{equation}\label{eq:finite-dimensional-operator-coefficient-convergence}
\lim_{j\to\infty}\alpha_{i,j}=0.
\end{equation}
于是,
\begin{equation}\label{eq:finite-dimensional-operator-continuity-proof}
\lim_{j\to\infty}A(x_j)=\lim_{j\to\infty}\sum_{i=1}^n\alpha_{i,j}A(e_i)=0.
\end{equation}

\end{Proof}
\endgroup

现在明确给出算子延拓的概念.

\begin{Definition}
\label{def:linear-operator-extension}
若 $\mathcal{D}(A)\subseteq\mathcal{D}(\widetilde A)$ 且
\begin{equation}\label{eq:linear-operator-extension-agreement}
A(x)=\widetilde A(x),\quad\text{for every }x\in\mathcal{D}(A).
\end{equation}
则称算子 $(\mathcal{D}(\widetilde A),\widetilde A)$ 为 $(\mathcal{D}(A),A)$ 的\textbf{延拓}.
\end{Definition}

以下定理断言, 可以将任意有界算子延拓到整个 $X$, 而不改变其范数.

\begin{Theorem}
\label{thm:bounded-operator-norm-preserving-extension}
每个有界算子 $(\mathcal{D}(A),A)$ 都可以在不改变其范数的情况下延拓到整个 $X$; 即存在延拓 $A$ 的有界算子 $(\mathcal{D}(\widetilde A),\widetilde A)$, 使 $\mathcal{D}(\widetilde A)=X$ 且 $\|\widetilde A\|=\|A\|$.
\end{Theorem}

\begingroup
\postdisplaypenalty=10000
\renewcommand{\endProof}{\par\nobreak\hfill$\blacksquare$\par\vspace{0.1cm}}
\begin{Proof}
只在两个特殊情形中证明这一点, 但它们包含大多数应用.

\textbf{Case 1.} \emph{定义域 $\mathcal{D}(A)$ 在 $X$ 中稠密.} 在此情形, 只需由连续性在整个 $X$ 上定义算子; 即对每个 $x\in X$, 存在序列 $x_n\in\mathcal{D}(A)$, 使 $x_n\to x$. 由于算子有界, 序列 $Ax_n$ 为 Cauchy 序列, 因而有极限. 于是可以定义\[Ax:=\lim_{n\to\infty}Ax_n.\]为说明这一定义无歧义, 注意对任意另一个收敛到 $x$ 的序列 $\widetilde x_n\in\mathcal{D}(A)$, 都有 $x_n-\widetilde x_n\to0$. 因而由连续性\[\lim_{n\to\infty}[Ax_n-A\widetilde x_n]=\lim_{n\to\infty}[A(x_n-\widetilde x_n)]=0.\]延拓后的算子具有相同范数, 因为\[x\mapsto\dfrac{\|Ax\|}{\|x\|}\]在 $x\ne0$ 处连续.

\textbf{Case 2.} \emph{$X=H$, 为 Hilbert 空间.} 若 $\mathcal{D}(A)$ 不闭, 就像上面的情形 1 一样, 由连续性将算子延拓到 $\overline{\mathcal{D}(A)}$. 由于 $H$ 为 Hilbert 空间, 现在可用投影定理, 通过下式定义满足 $\mathcal{D}(\widetilde A)=H$ 的算子 $(\mathcal{D}(\widetilde A),\widetilde A)$:
\begin{equation}\label{eq:hilbert-domain-operator-projection-extension}
\widetilde A(x)=\begin{cases}A(x),&x\in\overline{\mathcal{D}(A)}\\0,&x\in\overline{\mathcal{D}(A)}^{\perp}\end{cases}
\end{equation}
再由线性性延拓到 $H$ 的其余部分. $\|\widetilde A\|=\|A\|$ 的证明留给读者.

当 $X$ 为 Banach 空间时, 若 $\overline{\mathcal{D}(A)}\ne X$, 则缺少正交分解, 必须用一些巧妙技巧绕过这一困难. 由于这在应用中很少发生, 在证明中略去这一情形.
\end{Proof}
\endgroup

由于 Theorem~\ref{thm:bounded-operator-norm-preserving-extension}, 几乎总是假设有界算子的定义域为整个空间 $X$.

在 Chapter~\ref{chap:function-spaces} 中看到, 配备算子范数的 $\mathcal{L}(X,Y)$ 为 Banach 空间. 下文有时也关注逐点收敛, 而不是范数收敛. 作如下定义.

\begin{Definition}
\label{def:operator-strong-convergence}
若对每个 $x\in X$, 都有 $A_nx\to Ax$, 则称 $A_n\in\mathcal{L}(X,Y)$ \textbf{强收敛}到 $A\in\mathcal{L}(X,Y)$.
\end{Definition}

例如, 设 $X=Y=\ell^2$, $A_n$ 为在 $n$ 项之后截断序列的算子. 则 $A_n$ 强收敛到恒等算子, 但不在算子范数意义下收敛.
